\chapter*{Conclusion} 
\addcontentsline{toc}{chapter}{Conclusion} 
\markboth{Conclusion}{Conclusion} 

In the thesis, we focused on two main goals. Firstly, to design and build a machine-learning-based prediction model capable of predicting the next-day price
of individual stocks from the chosen market segment. Secondly, we focused on employing the acquired models in subsequent portfolio optimization 
and strategy building using the MVP framework.

Chapter \ref{chap:1} introduced the fundamental concepts of ML and modeling to the reader. The chapter presented the foundations of all algorithms and frameworks used in the thesis
and briefly overviewed the concepts of ML theory, which are frequently referred to in the thesis.

A specific training, validation, and testing splitting concept was adopted in the thesis. Taking into account the day-to-day prediction methodology, we adopted a rolling window, where 
the latest available predictor data (acquired from the constantly long interval) are used to predict the next-day price of a particular stock. On top of that, we used a unified set of predictors, consisting 
largely of the market data. Chapter \ref{chap:2} introduced the concepts of the rolling window, data preparation and the set of predictors used in the thesis.

Subsequently, it was possible to further evaluate the adopted ML frameworks and optimize their settings in order to achieve the best prediction results. The results showed that stocks in the IT segment,
in the particular time interval, were difficult to predict. Naïve prediction showed a great potential and justified that stock price development aligns more with the random walk \cite{fama1970efficient}. 

The introduced metrics of dominance and MAPE showed that vast majority of employed models were outperformed by the Naïve prediction in terms of prediction accuracy.
Only one framework (ARIMA model) outperformed the Naïve model in terms of dominance, which showed that time series models might be more capable of generalizing the pattern than conventional ML models.
This result is often not stated in a number of studies (e.g., \cite{patel}, \cite{vijh}, \cite{konferencia1}), which might lead to an overestimation of the models' practical utility.

Subsequent analyses of the application of the models to the portfolio optimization were introduced in Chapter \ref{chap:4}. Evaluation results suggested that even though prediction of most of the models were inferior to the Naïve benchmark,
the performance of portfolios, backed by models, was above-average compared to the conventional stock indices. Moreover, the portfolio backed by the Naïve model delivered comparable results to other non-naïve models.

On the other hand, quite a counterintuitive fact arose in the analyses. While during the period between 2017 and 2022, models with the lowest dominance delivered the best-performing portfolios, during the period between 2023 and 2026 the opposite
emerged. The portfolio backed by the Lasso model (which was evaluated as the least dominant model) was the best performer (on the testing set) in terms of average annual return across all the considered models. Even so, the NASDAQ index outperformed the portfolios, which is opposite
to the previous data sample as well.

Comparing both analyzed intervals, the results of the research suggest that the proposed portfolio models (accompanied with the prediction modeling) tend to be less volatile and profitable during periods of neither a bull nor a bear market than Naïve portfolio or conventional indices. On the other hand, during a strong bull market (in this case, the technology segment),
strategy performed worse than the NASDAQ index.

Aligning with Palomar \cite{portfolio}, the actual predictions tend to be less important for the portfolio optimization than the actual covariance of returns -- realized volatility of the stocks.
\newline
\par
Apart from the compared studies (e.g., \cite{krauss}, \cite{konferencia1}, \cite{zhao}), possibilities of applying such predictions for portfolio management were researched. These studies focused solely on designing the prediction pipeline and prediction evaluation. Connecting this methodology with portfolio optimization increases the application value of the research
and brings an unconventional point of view regarding the model's evaluation.

Used ML algorithms cannot comprehend the time-flow dependency properly \cite{forecast_book}. However, by using the ARIMA model we addressed the issue. Even though time series models are used in the research quite often, simpler frameworks are used more, even with inferior performance \cite{review}. This suggested a suitable alternative to the conventional ML frameworks. 

As stated earlier, prediction Naïve benchmark showed competitive results, in comparison with the other models. This shows that there exists a less computationally heavy prediction alternative, capable of outperforming advanced ML frameworks. Some studies tend to compare models between each other \cite{review} or omit the comparison at all (e.g., \cite{cijh}, \cite{zhao}). The thesis shows potential
worse performance of the models compared to the computationally inexpensive approaches.

Quantitative and qualitative comparison with the existing studies is limited. Direct comparison cannot be done, as there is no widely known study conducting similar rolling-window prediction on the whole IT market segment. However, several implicit conclusions can be done.
Sezer et al. \cite{review} suggest that many studies show various best performing frameworks. This suggests that particular models might perform diversely on various assets across the history.

Moreover, the ML workflow of applying models out-of-sample tends to be maintained \cite{krauss,review,zhao}. However, majority of studies used static time series splitting – static time points break the whole data segment into particular ML sets. Minority of studies implemented rolling window methodology \cite{review}. 
\newline
\par
Such an analysis of any coherent group of stocks could be performed using the designed pipeline. For future developments, it would be beneficial to conduct such research for the whole US market, which might 
create a trustworthy foundations for informed decision-making in the market. Lastly, by researching and implementing more sophisticated portfolio optimization frameworks, the overall performance of the investment 
might improve.

Apart from that, there is possibility of implementing more advanced frameworks – neural networks. Some studies suggest more reliable prediction results than conventional ML algorithms \cite{review}. However, we decided not to include neural networks in this research, as the data preparation implemented does not align with the requirements of other used frameworks. Altering the preparation process 
just for the purposes of neural networks would bias the analyses. 

Lastly, regarding the goals of the thesis, all of the targeted aims were fully accomplished.